\documentclass[11pt]{article}
\usepackage[utf8]{inputenc}
\usepackage{amsmath,amssymb}
\usepackage[margin=1in]{geometry}
\usepackage{hyperref}
\emergencystretch=3em

\title{The polynomial Taylor tree in $d$ dimensions}
\author{Claude, with Daniel Murfet}
\date{2026-09-06}

\begin{document}
\maketitle
\sloppy

\begin{abstract}
For a polynomial perturbation $J(u)=\sum_{\gamma\in S}c_\gamma u^\gamma$ the multinomial theorem
expands $(J(u))^p$ into monomials $u^{\sum_\gamma\kappa_\gamma\gamma}$ indexed by compositions
$\kappa\in\texttt{piAntidiag}\,S\,p$, so each layer of the exponential tree (unit~62) is a finite sum
of tree terms (unit~60) and the chart standard integral with $\xi=a+J$ is the exact double series
$\sum_p\frac{\beta^p}{p!}\sum_\kappa\binom{p}{\kappa}\prod_\gamma c_\gamma^{\kappa_\gamma}\cdot
\texttt{treeTerm}(\textstyle\sum_\gamma\kappa_\gamma\gamma,p)$: the paper's
\texttt{eq:flucttreeterms} grouped by composition. Zero \texttt{sorry}/\texttt{axiom}.
\end{abstract}

\section{Setting}
Mathlib's \texttt{Finset.sum\_pow\_eq\_sum\_piAntidiag} is the multinomial theorem over a finite index
set; the only bookkeeping is the identification
$\prod_\gamma(\prod_iu_i^{\gamma_i})^{\kappa_\gamma}=\prod_iu_i^{\sum_\gamma\kappa_\gamma\gamma_i}$,
which is \texttt{Finset.prod\_comm} plus \texttt{Finset.prod\_pow\_eq\_pow\_sum}.

\section{The things formalised}
{\small
\begin{verbatim}
noncomputable def polyJ {d : ℕ} (S : Finset (Fin d → ℕ)) (coef : (Fin d → ℕ) → ℝ)
    (u : Fin d → ℝ) : ℝ := ∑ γ ∈ S, coef γ * ∏ i, u i ^ γ i
def compShift {d : ℕ} (S : Finset (Fin d → ℕ)) (κ : (Fin d → ℕ) → ℕ) (i : Fin d) : ℕ :=
  ∑ γ ∈ S, κ γ * γ i
theorem polyJ_pow ... : polyJ S coef u ^ p = ∑ κ ∈ Finset.piAntidiag S p,
      (Nat.multinomial S κ : ℝ) * (∏ γ ∈ S, coef γ ^ κ γ) * ∏ i, u i ^ compShift S κ i
theorem treeLayer_polynomial ... : treeLayer β a b c k h (polyJ S coef) p
      = β ^ p / p.factorial * ∑ κ ∈ Finset.piAntidiag S p,
          (Nat.multinomial S κ : ℝ) * (∏ γ ∈ S, coef γ ^ κ γ)
            * treeTerm β a b c k h (compShift S κ) p
theorem boxIntegralXi_polynomial_hasSum ... :
    HasSum (fun p => β ^ p / p.factorial * ∑ κ ∈ Finset.piAntidiag S p, ...)
      (boxIntegralXi β a b c k h (polyJ S coef))
\end{verbatim}
}

\section{Proof strategy}
The pointwise expansion is the multinomial theorem followed by \texttt{prod\_mul\_distrib},
\texttt{prod\_comm}, \texttt{pow\_mul} and \texttt{prod\_pow\_eq\_pow\_sum}. The layer identity
distributes the finite sum through the integral (\texttt{integral\_finsetSum}, integrability from
continuity on the compact box) and pulls out the constants; the series statement is
\texttt{HasSum.congr\_fun} of unit~62.

\section{Roadblocks and resolutions}
\begin{itemize}
\item \texttt{Nat.multinomial} and the multinomial theorem live in
\texttt{Mathlib.Data.Nat.Choose.Multinomial}, outside the import closure.
\item \texttt{rw [Finset.mul\_sum, Finset.sum\_mul]} fails when the sum is nested under two products;
\texttt{simp only} with both lemmas distributes everything.
\item \texttt{integral\_finsetSum} inferred its summand from the integrability proof in Pi form; giving
\texttt{(f := …)} explicitly as a lambda makes the rewrite match.
\end{itemize}

\section{What was Mathlib, what was new}
Mathlib: \texttt{Finset.sum\_pow\_eq\_sum\_piAntidiag}, \texttt{Nat.multinomial},
\texttt{Finset.piAntidiag}, \texttt{Finset.prod\_comm}, \texttt{Finset.prod\_pow\_eq\_pow\_sum},
\texttt{integral\_finsetSum}. New: the polynomial tree expansion.

\section{Lessons}
The paper's grouping of the $p$-fold products of Taylor coefficients by total multi-index is exactly a
sum over compositions in Mathlib's \texttt{piAntidiag}; formalising along that grouping avoided any
multiset combinatorics.

\section{Outcome}
For polynomial $\xi$ the full Taylor tree of \texttt{thm:TaylorTree} is an exact, convergent double
series of tree terms, each of which has leading exponent in $\Lambda(h,k)$ and logarithmic degree at
most $d-1$. The analytic remainder of the theorem is the resummation of this series.
\end{document}
